\documentclass[12pt,a4paper]{article}
\usepackage[utf8]{inputenc} % inutile avec XeLaTeX/LuaLaTeX
\usepackage[T1]{fontenc}
\usepackage{amsmath,amssymb,mathrsfs,tikz,times,pifont}
\usepackage{enumitem}
\usepackage{multicol}
\usepackage{lmodern}
\newcommand\circitem[1]{%
\tikz[baseline=(char.base)]{
\node[circle,draw=gray, fill=red!55,
minimum size=1.2em,inner sep=0] (char) {#1};}}
\newcommand\boxitem[1]{%
\tikz[baseline=(char.base)]{
\node[fill=cyan,
minimum size=1.2em,inner sep=0] (char) {#1};}}
\setlist[enumerate,1]{label=\protect\circitem{\arabic*}}
\setlist[enumerate,2]{label=\protect\boxitem{\alph*}}
\everymath{\displaystyle}
\usepackage[left=1cm,right=1cm,top=1cm,bottom=1.7cm]{geometry}
\usepackage[colorlinks=true, linkcolor=blue, urlcolor=blue, citecolor=blue]{hyperref}
\usepackage{array,multirow}
\usepackage[most]{tcolorbox}
\usepackage{varwidth}
\usepackage{float}
\tcbuselibrary{skins,hooks}
\usetikzlibrary{patterns}

\newtcolorbox{exa}[2][]{enhanced,breakable,before skip=2mm,after skip=5mm,
colback=yellow!20!white,colframe=black!20!blue,boxrule=0.5mm,
attach boxed title to top left ={xshift=0.6cm,yshift*=1mm-\tcboxedtitleheight},
fonttitle=\bfseries,
title={#2},#1,
boxed title style={frame code={
\path[fill=tcbcolback!30!black]
([yshift=-1mm,xshift=-1mm]frame.north west)
arc[start angle=0,end angle=180,radius=1mm]
([yshift=-1mm,xshift=1mm]frame.north east)
arc[start angle=180,end angle=0,radius=1mm];
\path[left color=tcbcolback!60!black,right color = tcbcolback!60!black,
middle color = tcbcolback!80!black]
([xshift=-2mm]frame.north west) -- ([xshift=2mm]frame.north east)
[rounded corners=1mm]-- ([xshift=1mm,yshift=-1mm]frame.north east)
-- (frame.south east) -- (frame.south west)
-- ([xshift=-1mm,yshift=-1mm]frame.north west)
[sharp corners]-- cycle;
},interior engine=empty,
},interior style={top color=yellow!5}}

\usepackage{fancyhdr}
\usepackage{eso-pic}
\usepackage{tkz-tab}
\AddToShipoutPicture{
    \AtTextCenter{%
        \makebox[0pt]{\rotatebox{80}{\textcolor[gray]{0.7}{\fontsize{5cm}{5cm}\selectfont PGB}}}
    }
}

\usepackage{verbatim}

\usepackage{color,soul}

\usepackage{amsmath}
\usepackage{amsfonts}
\usepackage{amssymb}
\usepackage{systeme}
\usepackage{tkz-tab}
\usepackage{tikz}
\usetikzlibrary{arrows}
\newcounter{exemple} % Compteur pour les questions

% Définir la commande pour afficher une question numérotée
\newcommand{\exemple}{%
  \refstepcounter{exemple}%
  \textbf{\textcolor{green}{Exemple \theexemple :}} \ignorespaces
}
%---------------------------------------
\definecolor{myorange}{rgb}{1.0, 0.8, 0.0}

% Définir un compteur pour les exercices d'application
\newcounter{exerciceapp}

% Définir la commande pour afficher un exercice d'application numéroté
\newcommand{\exerciceapp}{%
  \refstepcounter{exerciceapp}%
  \textbf{\textcolor{myorange}{Exercice d'application \theexerciceapp :}} \ignorespaces
}
%--------------------------------------
% Définir un compteur pour les exercices d'application
\newcounter{correction}

% Définir la commande pour afficher un correction exercice d'application numéroté
\newcommand{\correction}{%
  \refstepcounter{correction}%
  \textbf{\textcolor{myorange}{Correction \thecorrection :}} \ignorespaces
}
%--------------------------------------
% Commande pour ajouter du texte en arrière-plan
\usepackage{fancyhdr}
\usepackage{eso-pic}
%\usepackage{tkz-tab}
\AddToShipoutPicture{
    \AtTextCenter{%
        \makebox[0pt]{\rotatebox{80}{\textcolor[gray]{0.7}{\fontsize{5cm}{5cm}\selectfont PGB}}}
    }
}
%This command takes a colour as an optional argument; the default colour is black.
\usetikzlibrary{shapes.geometric,fit}
\newcommand{\myul}[2][black]{\setulcolor{#1}\ul{#2}\setulcolor{black}}
\newcommand\tab[1][1cm]{\hspace*{#1}}

\begin{document}
% En-tête personnalisée
\begin{center}
    \Large\textbf{\underline{ EQUATIONS A UNE INCONNUE }}\\[-0.1cm]
    \normalsize\textbf{Prof : M. BA} \hfill \textbf{Classe : 4eme}\\[-0.1cm]
    \textbf{Année scolaire : 2025 -- 2026}
\end{center}



\subsection*{\textcolor{red}{\underline{Introduction :}}}

Dans la vie courante, nous rencontrons souvent des situations où une quantité inconnue doit être déterminée.  
Mettre une situation en équation consiste à traduire un problème concret en langage mathématique afin de trouver cette valeur inconnue.


\section*{\underline{\textcolor{red}{I - METTRE EN ÉQUATION UNE SITUATION SIMPLE}}}

\subsection*{\textcolor{red}{\underline{Objectif :}}}
\textbf{Savoir traduire une situation concrète par une équation et résoudre une équation du premier degré à une inconnue.}

\vspace{0.3cm}

\subsection*{\textcolor{red}{\underline{I - 1) Activité :}}}

Omar achète une gomme et un crayon noir pour 150 F.  
La gomme coûte deux fois plus cher que le crayon noir.

\begin{enumerate}
    \item On désigne par $x$ le prix du crayon noir.
    
    \begin{enumerate}
        \item Exprimer le prix de la gomme en fonction de $x$.
        \item Exprimer la dépense totale en fonction de $x$.
        \item Déterminer le prix du crayon noir.
    \end{enumerate}
\end{enumerate}

\vspace{0.5cm}

\textbf{Mise en équation :}

Si le prix du crayon noir est $x$, alors le prix de la gomme est $2x$.

La dépense totale est donc :
\[
x + 2x = 150
\]

On obtient l’équation :
\[
3x = 150
\]

\[
x = 50
\]

Le prix du crayon noir est donc \textbf{50 F} et celui de la gomme est \textbf{100 F}.

\vspace{0.5cm}

\subsection*{\textcolor{red}{\underline{I - 2) Définition :}}}

\textbf{Définition :}  
Une \textbf{équation} est une égalité contenant au moins une inconnue.  

L’inconnue est généralement désignée par une lettre telle que $x$, $y$, $a$, etc.

\vspace{0.3cm}

\textbf{Exemples :}

\[
x - 3 = 0 \quad ; \quad 7a - 5 = 2a + 3 \quad ; \quad 3m + 9 = 0 \quad ; \quad -4y + 12 = \frac{3x + 7}{3}
\]

\vspace{0.5cm}

\section*{\underline{\textcolor{red}{II - RÉSOLUTION D’UNE ÉQUATION}}}

Résoudre une équation dans $\mathbb{Q}$ revient à déterminer l’ensemble des nombres rationnels qui vérifient cette égalité.  

Cet ensemble est appelé \textbf{ensemble des solutions} de l’équation.  
On le note $S$.

\vspace{0.4cm}

\underline{\exemple}

\begin{enumerate}
\item Résoudre dans $\mathbb{Q}$ l’équation $x + 6 = 9$.

le nombre qu'il faut ajouter à 6 pour avoir 9 est 3
\[
x = 3
\]

Comme $3 \in \mathbb{Q}$, l’ensemble des solutions est :
\[
S = \{3\}
\]

\vspace{0.5cm}

\item Résoudre dans $\mathbb{Q}$ l’équation $3x = 150$.

le nombre qu'il faut multiplier par 3 pour avoir 150 est 50

\[
x = 50
\]

Comme $50 \in \mathbb{Q}$, l’ensemble des solutions est :
\[
S = \{50\}
\]
\end{enumerate}

\subsection*{\underline{\textcolor{red}{A - Équations se ramenant à la forme $ax + b = 0$ avec $a \neq 0$}}}

\vspace{0.3cm}

\subsubsection*{\textcolor{red}{1) Définition}}

On appelle équation du premier degré à une inconnue toute équation pouvant s’écrire sous la forme :

\[
ax + b = 0
\]

où $a$ et $b$ sont des nombres rationnels et $a \neq 0$.

\vspace{0.4cm}

\subsubsection*{\textcolor{red}{2) Méthode de résolution}}

Pour résoudre une équation du type $ax + b = 0$, il faut :

\begin{itemize}
    \item[\textcolor{red}{\checkmark}] ajouter à chaque membre de l’égalité l’opposé de $b$ ;
    \item[\textcolor{red}{\checkmark}] multiplier chaque membre de la nouvelle équation par l’inverse de $a$ ;
    \item[\textcolor{red}{\checkmark}] vérifier, puis donner la solution de l’équation.
\end{itemize}

\vspace{0.4cm}

\textbf{Soit à résoudre l’équation $ax + b = 0$ avec $a \neq 0$.}

\[
ax + b = 0
\]
\textit{On part de l’équation donnée.}

\[
ax + b - b = 0 - b
\]
\textit{On ajoute $-b$ aux deux membres de l’égalité afin de supprimer le terme $b$ dans le premier membre.}

\[
ax = -b
\]
\textit{Après simplification, il reste uniquement le terme contenant l’inconnue.}

\[
x = \frac{-b}{a}
\]
\textit{On multiplie les deux membres par l’inverse de $a$ (possible car $a \neq 0$) afin d’isoler $x$.}

d’où 
\[
S = \left\{ \frac{-b}{a} \right\}
\]
\textit{L’équation admet donc une solution unique.}

\vspace{0.5cm}

\subsubsection*{\textcolor{red}{3) Exemples}}

\subsubsection*{\textcolor{red}{Exemple :}}

Résoudre dans $\mathbb{Q}$ l’équation :

\[
5x - 15 = 0
\]

\[
5x - 15 + 15 = 0 + 15
\]
\textit{On ajoute $15$ aux deux membres pour supprimer le terme constant.}

\[
5x = 15
\]
\textit{Après simplification, il reste le terme contenant l’inconnue.}

\[
x = \frac{15}{5}
\]
\textit{On divise les deux membres par $5$ (possible car $5 \neq 0$) afin d’isoler $x$.}

\[
x = 3
\]

Donc :
\[
S = \{3\}
\]
\textit{L’équation admet une solution unique.}

\subsubsection*{\textcolor{red}{Exemple :}}

Résoudre dans $\mathbb{Q}$ l’équation :

\[
-4x + 12 = 0
\]

\[
-4x + 12 - 12 = 0 - 12
\]
\textit{On ajoute $-12$ aux deux membres pour supprimer le terme constant.}

\[
-4x = -12
\]
\textit{Après simplification, il reste le terme contenant l’inconnue.}

\[
x = \frac{-12}{-4}
\]
\textit{On divise les deux membres par $-4$ (possible car $-4 \neq 0$) afin d’isoler $x$.}

\[
x = 3
\]

Donc :
\[
S = \{3\}
\]
\textit{L’équation admet une solution unique.}

\subsubsection*{\textcolor{red}{4) Remarque importante}}

Dans la pratique, au lieu de dire que l’on ajoute l’opposé de $b$ aux deux membres, 
on dit souvent que l’on \textbf{transpose} $b$ dans l’autre membre.

Ainsi, dans l’équation :

\[
ax + b = 0
\]

on écrit directement :

\[
ax = -b
\]

\textbf{Règle 1 :} Lorsqu’un \textbf{terme} change de membre, il change de signe.

\medskip

\begin{comment}
En réalité, transposer un terme consiste à ajouter son opposé aux deux membres de l’égalité.

Par exemple :

\[
ax - b = 0 \quad \Longrightarrow \quad ax = b
\]

\textit{Le terme $-b$ devient $+b$.}
\end{comment}

\vspace{0.4cm}

\textbf{Attention :}

Si c’est un \textbf{facteur} qui change de membre, il \textbf{ne change pas de signe}.  
On effectue alors une division (ou une multiplication).

Par exemple :

\[
3x = 12
\]

\[
x = \frac{12}{3}
\]

\textit{Le nombre $3$ ne change pas de signe ; on divise simplement par $3$.}


\textbf{\exerciceapp}

Résoudre dans $\mathbb{Q}$ les équations suivantes :

\begin{enumerate}
    \item $7x + 5 = 0$
    
    \item $-3x - 9 = 0$
    
    \item $4x = 20$
    
    \item $-5x + 15 = 0$
    
    \item $6 - 2x = 0$
\end{enumerate}
\textbf{\correction}\\

\subsection*{\underline{\textcolor{red}{B - Passage de $cx + d = ex + f$ à la forme $ax + b = 0$ }}}

\subsection*{\textcolor{red}{1) Principe}}

On considère une équation du type :

\[
cx + d = ex + f
\]

où $c$, $d$, $e$ et $f$ sont des nombres rationnels.

L’objectif est de regrouper tous les termes dans un même membre afin d’obtenir une équation de la forme :

\[
ax + b = 0
\]

\vspace{0.5cm}

\subsection*{\textcolor{red}{2) Méthode}}

\[
cx + d = ex + f
\]

\textbf{Étape 1 :} On regroupe les termes contenant $x$ dans le premier membre.

\[
cx - ex + d = f
\]

\[
(c - e)x + d = f
\]

\textbf{Étape 2 :} On regroupe les termes constants dans le second membre.

\[
(c - e)x = f - d
\]

\textbf{Étape 3 :} On met tout dans le même membre.

\[
(c - e)x - (f - d) = 0
\]

On obtient alors une équation de la forme :

\[
ax + b = 0
\]

avec :
\[
a = c - e \quad \text{et} \quad b = -(f - d)
\]

\vspace{0.5cm}

\subsection*{\textcolor{red}{3) Résolution}}

Si $a \neq 0$, l’équation admet une solution unique :

\[
x = \frac{-b}{a}
\]

\vspace{0.5cm}

\subsection*{\textcolor{red}{4) Exemple}}

Résoudre dans $\mathbb{Q}$ l’équation :

\[
3x + 5 = x + 9
\]

\[
3x - x + 5 = 9
\]

\[
2x + 5 = 9
\]

\[
2x = 9 - 5
\]

\[
2x = 4
\]

\[
x = 2
\]

Donc :
\[
S = \{2\}
\]

\textbf{\exerciceapp}

Résoudre dans $\mathbb{Q}$ les équations suivantes en les ramenant d’abord à la forme $ax + b = 0$ :

\begin{enumerate}
    \item $4x + 3 = 2x + 11$
    
    \item $7x - 5 = 3x + 9$
    
    \item $5x + 8 = 5x - 2$
    
    \item $6x - 4 = 6x + 1$
    
    \item $-2x + 7 = 3x - 8$
\end{enumerate}

\section*{\underline{\textcolor{red}{III - Équations de la forme $(ax + b)(cx + d) = 0$}}}

\subsection*{\textcolor{red}{1) Principe fondamental}}

On considère une équation de la forme :

\[
(ax + b)(cx + d) = 0
\]

où $a$, $b$, $c$ et $d$ sont des nombres rationnels.

\medskip

\textbf{Propriété (Produit nul) :}

Un produit est nul si et seulement si au moins un de ses facteurs est nul.

Autrement dit :

\[
A \times B = 0 \quad \Longleftrightarrow \quad A = 0 \ \text{ou} \ B = 0
\]

\vspace{0.5cm}

\subsection*{\textcolor{red}{2) Méthode de résolution}}

Pour résoudre :

\[
(ax + b)(cx + d) = 0
\]

On résout séparément les deux équations :

\[
ax + b = 0
\quad \text{ou} \quad
cx + d = 0
\]

On obtient :

\[
x = \frac{-b}{a}
\quad \text{ou} \quad
x = \frac{-d}{c}
\]

L’ensemble des solutions est donc :

\[
S = \left\{ \frac{-b}{a} \ ; \ \frac{-d}{c} \right\}
\]

(à condition que $a \neq 0$ et $c \neq 0$)

\vspace{0.5cm}

\subsection*{\textcolor{red}{3) Exemple}}

Résoudre dans $\mathbb{Q}$ l’équation :

\[
(2x - 4)(x + 3) = 0
\]

D’après la propriété du produit nul :

\[
2x - 4 = 0
\quad \text{ou} \quad
x + 3 = 0
\]

Première équation :

\[
2x = 4
\]

\[
x = 2
\]

Deuxième équation :

\[
x = -3
\]

Donc :

\[
S = \{2 \ ; \ -3\}
\]
\textbf{\exerciceapp}

\begin{enumerate}
\item Résoudre dans $\mathbb{Q}$ les équations ci-dessous :
\begin{enumerate}
    \item $(3x - 6)(x + 4) = 0$
    
    \item $(2x + 5)(5x - 10) = 0$
    
    \item $(4x - 8)(2x + 3) = 0$
    
    \item $(x - 7)(-3x + 9) = 0$
    
    \item $(5x + 1)(5x + 1) = 0$
\end{enumerate}

\item Résoudre dans $\mathbb{Q}$ les équations ci-dessous :
\begin{enumerate}
\item $x^2 - 9 = 0 $
\item $4x^2 - 25 = 0$
\item $(2x - 3)^2 = (2x - 3)(x + 4)$
\item $4x^2 - 1 + (2x - 1)(x + 3) = 0$
\end{enumerate}
\end{enumerate}
\section*{\underline{\textcolor{red}{IV - Équations quotient}}}
\subsection*{\underline{\textcolor{red}{A. - Équations quotient du type $\dfrac{P(x)}{Q(x)} = 0$}}}
\subsection*{\underline{\textcolor{red}{1) Définition}}}

On appelle \textbf{équation quotient} toute équation pouvant s’écrire sous la forme :

\[
\frac{p(x)}{q(x)} = 0
\]

où $p(x)$ et $q(x)$ sont des expressions algébriques et $q(x)$ n’est pas identiquement nul.

\medskip

L’équation n’a de sens que pour les valeurs de $x$ telles que $q(x) \neq 0$.

\vspace{0.5cm}

\subsection*{\underline{\textcolor{red}{2) Propriété fondamentale}}}

Un quotient est nul si et seulement si son numérateur est nul et son dénominateur est non nul.

Ainsi,

\[
\frac{p(x)}{q(x)} = 0
\quad \Longleftrightarrow \quad
\begin{cases}
p(x) = 0 \\
q(x) \neq 0
\end{cases}
\]

\vspace{0.5cm}

\subsection*{\underline{\textcolor{red}{3) Méthode de résolution}}}

Pour résoudre une équation quotient :

\begin{enumerate}
    \item Déterminer les valeurs qui annulent le dénominateur  
    (ce sont les \textbf{valeurs interdites}).
    
    \item Résoudre l’équation obtenue en annulant le numérateur.
    
    \item Vérifier que les solutions trouvées ne font pas annuler le dénominateur.
\end{enumerate}

\vspace{0.5cm}

\subsection*{\underline{\textcolor{red}{4) Exemple}}}

Résoudre dans $\mathbb{Q}$ l’équation :

\[
\frac{x^2 - 4}{x - 2} = 0
\]

\textbf{Condition d’existence :}

\[
x - 2 \neq 0 \quad \Longrightarrow \quad x \neq 2
\]

\textbf{Résolution :}

\[
x^2 - 4 = 0
\]

\[
(x - 2)(x + 2) = 0
\]

\[
x = 2 \quad \text{ou} \quad x = -2
\]

Or $x \neq 2$, donc la solution admise est :

\[
S = \{-2\}
\]
\textbf{\exerciceapp}

Résoudre dans $\mathbb{Q}$ les équations suivantes :

\begin{enumerate}
    \item $\displaystyle \frac{x - 3}{x + 2} = 0$
    
    \item $\displaystyle \frac{2x + 5}{x - 1} = 0$
    
    \item $\displaystyle \frac{x^2 - 9}{x - 3} = 0$
    
    \item $\displaystyle \frac{(x - 4)(x + 1)}{x - 4} = 0$
    
    \item $\displaystyle \frac{x^2 - 4x}{x} = 0$
\end{enumerate}
\subsection*{\underline{\textcolor{red}{B- Équations quotient du type $\dfrac{a}{x} =\dfrac{b}{c}$ avec $x\neq 0$ et $c\neq 0$}}}

\subsubsection*{\textcolor{red}{1) Principe}}

On considère une équation de la forme :

\[
\frac{a}{x} = \frac{b}{c}
\]

où $a$, $b$ et $c$ sont des nombres rationnels avec $c \neq 0$.

L’équation n’a de sens que pour $x \neq 0$.

\vspace{0.4cm}
\subsubsection*{\textcolor{red}{2) Méthode de résolution}}

On considère l’équation :

\[
\frac{a}{x} = \frac{b}{c}
\quad \text{avec } x \neq 0 \text{ et } c \neq 0.
\]

On multiplie les deux membres par le produit $xc$ (possible car $x \neq 0$ et $c \neq 0$) :

\[
xc \times \frac{a}{x} = xc \times \frac{b}{c}
\]

\[
ac = bx
\]

On isole alors $x$ :
$
x = \frac{ac}{b}
\quad \text{(si } b \neq 0\text{)}.
$

\subsubsection*{\textcolor{red}{ Remarque (\textit{produit en croix})}}  

\[
\frac{a}{x} = \frac{b}{c},
\]

le produit des moyens est égal au produit des extrêmes.

Autrement dit : $a \times c = b \times x.$

On isole alors $x$ :
$
x = \frac{ac}{b}
\quad \text{(si } b \neq 0\text{)}.
$

\subsubsection*{\textcolor{red}{3) Cas particulier}}

\begin{itemize}
    \item Si $b = 0$, l’équation devient :
    \[
    \frac{a}{x} = 0
    \]
    Or un quotient est nul si et seulement si son numérateur est nul.
    
    Donc :
    \[
    a = 0
    \]
    
    \begin{itemize}
        \item Si $a = 0$, tout $x \neq 0$ est solution.
        \item Si $a \neq 0$, l’équation n’admet aucune solution.
    \end{itemize}
\end{itemize}

\vspace{0.5cm}

\subsubsection*{\textcolor{red}{4) Exemple}}

Résoudre dans $\mathbb{Q}$ l’équation :

\[
\frac{6}{x} = \frac{3}{4}
\]

On multiplie les deux membres par $4x$ :

\[
4x \times \frac{6}{x} = 4x \times \frac{3}{4}
\]

\[
24 = 3x
\]

\[
x = 8
\]

Comme $8 \neq 0$, la solution est :

\[
S = \{8\}
\]
\end{document}